			\TemplateSection[K0E1A0E020E990E0207020E910E990E8A0E9F0E0A0E990E9F0E8A0E210101120202020212010100]{Data Structure}
				%\subsection{非递归线段树}
					% \subsubsection{区间加，区间求和}
					% 	\inputminted{cpp}{src/DataStructure/非递归线段树求和.cpp}
					%\subsubsection{区间加，区间求最大值}
						%\inputminted{cpp}{src/DataStructure/非递归线段树求最大值.cpp}
				\TemplateSubsection[KC58813C4C043C0EB10C1561CC5B63A0101010100]{回滚并查集}
					\TemplateCode{cpp}{src/DataStructure/回滚并查集.cpp}
				\TemplateSubsection[KC9117FC35C0DC50F13C5B63A0101010100]{莫队合集}
					\TemplateSubsubsection{普通莫队}
						\TemplateCode{cpp}{src/DataStructure/莫队.cpp}
					\TemplateSubsubsection{树上莫队}
						\TemplateCode{cpp}{src/DataStructure/树上莫队.cpp}
					\TemplateSubsubsection{带修莫队}
						单点修改, 区间查询颜色个数模板题, 块长为 $n^{2/3}$, 时间复杂度为 $n^{3/5}$
						\TemplateCode{cpp}{src/DataStructure/带修莫队.cpp}
					\TemplateSubsubsection{回滚莫队}
						\TemplateNote{src/DataStructure/回滚莫队.tex}	
						\TemplateCode{cpp}{src/DataStructure/回滚莫队.cpp}
					\TemplateSubsubsection{二次离线莫队}
						\TemplateNote{src/DataStructure/二次离线莫队.tex}
						\TemplateCode{cpp}{src/DataStructure/二次离线莫队.cpp}
				\TemplateSubsection[KCCF440D18D16CCF4C4D1D1190101010100]{树状数组}
					\TemplateCode{cpp}{src/DataStructure/bit.cpp}
				\TemplateSubsection[KC2AE3AC2F419CF5422C42016CF062BC3530DCCF4400101010100]{单点修改线段树}
					\TemplateCode{cpp}{src/DataStructure/单点修改线段树.cpp}
				\TemplateSubsection[KC73152C0BF13C5BE19CF062BC3530DCCF4400101010100]{懒标记线段树}
					\TemplateCode{cpp}{src/DataStructure/懒标记线段树.cpp}
				\TemplateSubsection[KCF062BC3530DCCF440C50F13C0EB10D0470DC3CF0DC7BBCA0101010100]{线段树合并与分裂}
					\TemplateCode{cpp}{src/DataStructure/线段树合并与分裂.cpp}
				\TemplateSubsection[KC2EA8E07020E360702C29E0DC77B1CC1852ECF062BC3530DCCF4400101010100]{第 k 大李超线段树}
					\TemplateCode{cpp}{src/DataStructure/kthLiChao.cpp}
				\TemplateSubsection[KC2E343C6740DC39716CCF4400101010100]{笛卡尔树}
					% \input{src/DataStructure/tree.tex}
					\TemplateCode{cpp}{src/DataStructure/笛卡尔树.cpp}
				\TemplateSubsection[KC2F419C3CF0DD141760101010100]{点分治}
					\TemplateCode{cpp}{src/DataStructure/点分治.cpp}
				\TemplateSubsection[KC2F419C3CF0DCCF4400101010100]{点分树}
					\TemplateNote{src/DataStructure/点分树.tex}
				\TemplateSubsection[KCCF440CC9613CB0549C39D0DCCD931C50F13C0EB100101010100]{树上启发式合并}
					\TemplateCode{cpp}{src/DataStructure/dsu_on_tree.cpp}
				\TemplateSubsection[K0E7C07270E7007020E70072A070207020E7C07270D1A072A07020E480E0A0E020101120202020202020212020202021212120101FFF8821200]{$O(n \log n) - O(1)\ \text{ LCA}$}
					\TemplateCode{cpp}{src/DataStructure/o1lca.cpp}
				% \subsection{KD 树}
				% 	\inputminted{cpp}{src/DataStructure/KD tree.cpp}
				\TemplateSubsection[K0E910E7E0E480E020EA7010112010100]{Splay}
					\TemplateCode{cpp}{src/DataStructure/Splay.cpp}
				\TemplateSubsection[K0E480E0A0E990702C32E0DCD9F2BCCF4400101121212010100]{LCT 动态树}
					\TemplateCode{cpp}{src/DataStructure/LCT.cpp}
				\TemplateSubsection[KC6A50DC1BE46C63110C55F1007020E990E8A0E210E020E7E0101020202020212010100]{可持久化 Treap}
					\TemplateCode{cpp}{src/DataStructure/可持久化treap-zmc.cpp}
				\TemplateSubsection[K0E7E0E090E1A0E910702CCD255CF0622CAD113C52A70CCF440010112121212010100]{PBDS 实现平衡树}
					\TemplateCode{cpp}{src/DataStructure/pbds.cpp}
				\TemplateSubsection[K0E7E0E090E1A0E910702CCD255CF0622C6A50DC0EB10C3562E010112121212010100]{PBDS 实现可并堆}
					\TemplateCode{cpp}{src/DataStructure/pbds_heap.cpp}
				\TemplateSubsection[KCCE40DCF2B0D07020E090E320E990E910E210E990101010100]{手写 bitset}
					\TemplateCode{cpp}{src/DataStructure/bitset.cpp}
				\TemplateSubsection[KC69E16C36E0DC77ED3CCF4400101010100]{珂朵莉树}
					\TemplateCode{cpp}{src/DataStructure/odt.cpp}
				\TemplateSubsection[KD12F55CDDF0DC39A0DC3CF0D0101010100]{整体二分}
					\TemplateNote{src/DataStructure/binary.tex}
					\TemplateCode{cpp}{src/DataStructure/binary.cpp}
				\TemplateSubsection[KC32E0DCD9F2B0E1A0E7E0101010100]{动态dp}
					\TemplateNote{src/DataStructure/ddp.tex}
					\TemplateCode{cpp}{src/DataStructure/ddp.cpp}
				\TemplateSubsection[KC5B52ECBFE13CFE60DCF062BC3530DCCF4400101010100]{吉如一线段树}
					\TemplateNote{src/DataStructure/beats.tex}
					\TemplateCode{cpp}{src/DataStructure/beats.cpp}
				% \subsection{有旋 Treap}
				% 	\inputminted{cpp}{src/DataStructure/有旋treap.cpp}
					\newpage
